\documentclass[aspectratio=169,12pt]{beamer}
\usepackage{amsmath}
\usepackage{amssymb}
\usepackage{array}
\usepackage[most]{tcolorbox}
\usepackage{graphicx}
\usepackage[sfdefault,lf]{FiraSans}
\renewcommand{\familydefault}{\sfdefault}
\setlength{\parindent}{0pt}
\everymath{\displaystyle}

\setbeamertemplate{navigation symbols}{}
\setbeamertemplate{footline}{}
\setbeamertemplate{headline}{}
\setbeamertemplate{blocks}[rounded][shadow=false]

\definecolor{mmpageWhite}{HTML}{FCFBF7}
\definecolor{mmtanSoft}{HTML}{F7F5EF}
\definecolor{mmgreenDeep}{HTML}{244B2C}
\definecolor{mmgreenLeaf}{HTML}{5D8A56}
\definecolor{mmgreenSoft}{HTML}{E4EFD9}
\definecolor{mmtext}{HTML}{163221}
\definecolor{mmwarn}{HTML}{8F2F36}
\definecolor{mmwarnSoft}{HTML}{F8E8EA}

\setbeamercolor{background canvas}{bg=mmpageWhite}
\setbeamercolor{normal text}{fg=mmtext,bg=mmpageWhite}
\setbeamercolor{itemize item}{fg=mmgreenDeep}
\setbeamercolor{itemize subitem}{fg=mmgreenDeep}

\newtcolorbox{MMCard}[2][]{enhanced,breakable=false,colback=#2,colframe=black,boxrule=1.5pt,arc=8pt,left=5pt,right=5pt,top=4pt,bottom=4pt,before skip=0pt,after skip=0pt,#1}
\newcommand{\KoruIcon}{\raisebox{-0.2em}{\includegraphics[height=1.05em]{../../../manamaths/SITE/header-logo.png}}}
\newcommand{\NotesTitle}[1]{{\colorbox{mmtanSoft}{\parbox{0.975\linewidth}{\vspace{0.14em}\hspace{0.25em}{\LARGE\bfseries\textcolor{mmgreenDeep}{#1}\hfill\KoruIcon}\vspace{0.14em}}}\par\vspace{0.18em}{\color{mmgreenLeaf}\rule{\linewidth}{2.0pt}}\par\vspace{0.12em}}}
\newcommand{\SectionTitle}[1]{{\large\bfseries\textcolor{mmgreenDeep}{#1}}\par\vspace{0.04em}}
\newcommand{\GreenDivider}{\par\vspace{0.01em}{\color{mmgreenLeaf}\rule{\linewidth}{2.4pt}}\par\vspace{0.03em}}
\newcommand{\KeyIdea}[1]{\begin{MMCard}[title={\textbf{Key idea}},colbacktitle=mmgreenSoft,coltitle=mmgreenDeep]{mmtanSoft}\small #1\end{MMCard}}
\newcommand{\WorkedExample}[2]{\begin{MMCard}[title={\textbf{#1}},colbacktitle=mmgreenSoft,coltitle=mmgreenDeep]{white}\small #2\end{MMCard}}
\newcommand{\CommonMistake}[1]{\begin{MMCard}[title={\textbf{Common mistake}},colbacktitle=mmwarnSoft,coltitle=mmwarn]{white}\small #1\end{MMCard}}
\newcommand{\TryThese}[1]{\begin{MMCard}[title={\textbf{Try these}},colbacktitle=mmgreenSoft,coltitle=mmgreenDeep]{mmtanSoft}\small #1\end{MMCard}}
\newcommand{\StepLine}[2]{\textbf{#1.} #2\par\vspace{0.03em}}

\usepackage{tikz}
\setbeamertemplate{background}{
 \begin{tikzpicture}[remember picture,overlay]
 \draw[black,line width=1.5pt,rounded corners=10pt]
 ([xshift=0.45cm,yshift=-0.45cm]current page.north west)
 rectangle
 ([xshift=-0.45cm,yshift=0.45cm]current page.south east);
 \end{tikzpicture}
}

\begin{document}

\begin{frame}[t,shrink=10]
\NotesTitle{Notes \& Steps}
\footnotesize
\begin{columns}[T,onlytextwidth]
\begin{column}{0.48\textwidth}
\KeyIdea{Expanding a bracket means multiplying each term inside the bracket by the term outside. This uses the distributive property: \(a(b+c)=ab+ac\). The number (or variable) outside multiplies \textbf{every} term inside, one at a time.}
\SectionTitle{Steps}
\StepLine{1}{Identify the term outside the bracket (the multiplier).}
\StepLine{2}{Multiply it by the first term inside the bracket.}
\StepLine{3}{Multiply it by the next term inside the bracket.}
\StepLine{4}{Write as a sum: \(\text{term}_1 + \text{term}_2\).}
\StepLine{5}{Watch the signs! A negative outside flips the signs inside: \(-3(x-2) = -3x + 6\).}
\end{column}
\begin{column}{0.48\textwidth}
\SectionTitle{Examples}
\begin{itemize}
  \item \(4(x+3) = 4x + 12\)
  \item \(-2(y+5) = -2y - 10\)
  \item \(3(2x-1) = 6x - 3\)
  \item \(-5(4a-2) = -20a + 10\)
  \item \(x(x+4) = x^2 + 4x\)
  \item \(2a(3a-1) = 6a^2 - 2a\)
\end{itemize}
\end{column}
\end{columns}
\GreenDivider
\CommonMistake{Only multiplying the first term inside the bracket. The multiplier goes to \textbf{every} term. \(4(x+3) = 4x+12\), not \(4x+3\). Check by substituting \(x=1\): \(4(1+3)=16\) vs \(4(1)+3=7\) — only the first answer is right.}
\end{frame}

\begin{frame}[t,shrink=12]
\NotesTitle{Notes \& Steps}
\begin{columns}[T,onlytextwidth]
\begin{column}{0.48\textwidth}
\WorkedExample{Example 1: positive multiplier}{Expand \(5(2x+3)\).\[
5 \times 2x = 10x,\quad 5 \times 3 = 15
\]Answer: \(10x + 15\)}
\end{column}
\begin{column}{0.48\textwidth}
\WorkedExample{Example 2: negative multiplier}{Expand \(-4(3a-2)\).\[
-4 \times 3a = -12a,\quad -4 \times (-2) = 8
\]Answer: \(-12a + 8\)}
\end{column}
\end{columns}
\GreenDivider
\begin{columns}[T,onlytextwidth]
\begin{column}{0.48\textwidth}
\WorkedExample{Example 3: expand and simplify}{Simplify \(3(2x+4) - 5\).\[
3 \times 2x = 6x,\quad 3 \times 4 = 12
\]So \(6x + 12 - 5 = 6x + 7\)}
\end{column}
\begin{column}{0.48\textwidth}
\WorkedExample{Example 4: double expansion}{Simplify \(2(3x+1) + 4(2x-3)\).\[
6x+2 + 8x-12 = 14x - 10\]}
\end{column}
\end{columns}
\GreenDivider
\begin{columns}[b,onlytextwidth]
\begin{column}{0.48\textwidth}
\TryThese{\begin{enumerate}
  \item Expand \(3(x+7)\).
  \item Expand \(-2(4a-3)\).
  \item Expand and simplify \(4(2x+5)-3\).
\end{enumerate}}
\end{column}
\begin{column}{0.48\textwidth}
\CommonMistake{Forgetting to multiply the sign too. \(-3(x-4)\) gives \(-3x+12\), because \(-3 \times (-4) = +12\). Some students write \(-3x-12\) — that's wrong.}
\end{column}
\end{columns}
\end{frame}

\end{document}
